\documentclass[11pt]{article}
\usepackage{amsmath, amssymb, amsthm}
\usepackage{mathrsfs}
\usepackage{fullpage}
\usepackage{hyperref}

\title{\textbf{Language Specification: The Unified Generative Framework}\\
\large TOGT + GCM + dm$^3_{\mathrm{disc}}$\

\[4pt]
\large \emph{Discrete Generative Contact Mechanics (dm$^3_{\mathrm{disc}}$)}\\
\large A Formal Lift of the dm$^3$ Category to Arithmetic Dynamical Systems}
\author{Pablo Nogueira Grossi (Independent Researcher)}
\date{v1.0 --- April 2026}

\begin{document}
\maketitle

\begin{abstract}
This note provides the precise ``language-first'' specification requested by the
Collatz supplement (TOGT Report, pp.~108--109). It synthesizes:
\begin{itemize}
  \item the continuous dm$^3$ framework of the GCM Manifesto (\S2.2),
  \item the Topographical Orthogonal Generative Theory (TOGT) operator algebra
        $G = U \circ F \circ K \circ C$,
  \item the discrete lift dm$^3_{\mathrm{disc}}$ formalized in Lean~4
        (\texttt{DiscreteDm3.lean} v1.5).
\end{itemize}
The result is a single higher-order category in which both smooth contact flows
and piecewise-linear arithmetic maps (such as the Collatz map) are objects.
The Collatz conjecture is thereby reduced to a pure membership problem inside
this category.
\end{abstract}

\section{Continuous dm$^3$ (GCM Recap)}

A continuous dm$^3$-system is a tuple


\[
(M, X, \alpha, V, \Phi, \Gamma)
\]


where:
\begin{itemize}
  \item $M$ is a Riemannian manifold,
  \item $X$ is a $C^2$ vector field generating a flow,
  \item $\alpha$ is a contact form (maximally non-integrable hyperplane field),
  \item $V$ is a Lyapunov function with average descent outside the limit cycle,
  \item $\Phi$ is the contact Hamiltonian,
  \item $\Gamma$ is a hyperbolic structured limit cycle.
\end{itemize}

Morphisms are contactomorphisms $f : M \to M'$ such that:


\[
f \circ \varphi^t = \varphi'^t \circ f,\qquad
|\Phi'(f(x)) - \Phi(x)| \le C_\Phi,\qquad
V'(f(x)) \le V(x) + C_V.
\]



The GCM Manifesto (\S2.2) establishes four main theorems:
\begin{enumerate}
  \item existence and stability of $\Gamma$,
  \item categorical closure of the generative pipeline,
  \item contact normal form,
  \item structural stability under contact-preserving perturbations.
\end{enumerate}

\section{Discrete Lift: dm$^3_{\mathrm{disc}}$}

A discrete dm$^3$-system is a tuple


\[
(X, T, V, \Phi, \Gamma, \mathcal{G})
\]


exactly as in \texttt{DiscreteDm3System} (Lean v1.5), satisfying the axioms below.

\subsection{State Space}



\[
X = \mathbb{N}_{\ge 1}
\]


(or any countable discrete space) equipped with the 2-adic metric


\[
d_2(m,n) = 2^{-v_2(m-n)}.
\]



\subsection{Generative Map}

A piecewise-linear map $T : X \to X$.

\subsection{Discrete Contact Form}

\textbf{Axiom (contactForm).}  
For all odd $n$,


\[
v_2(3n+1) \ge 1.
\]


This is \emph{proved} for the Collatz map in Lean v1.5.

\subsection{Discrete Contact Hamiltonian}

A function $\Phi : X \to \mathbb{R}_{\ge 0}$.

\subsection{Lyapunov Function}

A function $V : X \to \mathbb{R}_{\ge 0}$.

\subsection{Structured Limit Cycle}

A finite attracting cycle $\Gamma$.

\subsection{Mean Contraction}

\textbf{Axiom (meanContraction).}  
There exist $\kappa < 1$ and $N_0 \in \mathbb{N}$ (the TOGT monster threshold $g_6 = 33$) such that


\[
\forall n \ge N_0,\qquad
\log(T^2(n)) - \log n \le \log \kappa.
\]


This encodes the classical $3/4$ macro-step contraction heuristic.

\subsection{Operator Grammar (TOGT)}

\textbf{Axiom (operatorDecomposition).}  


\[
T = G = U \circ F \circ K \circ C.
\]


This is \emph{proved} for the Collatz map in Lean v1.5.

\subsection{Categorical Closure}

The system belongs to the category dm$^3_{\mathrm{disc}}$.

\subsection{Morphisms}

Morphisms $f : A.X \to B.X$ satisfy:


\[
f \circ T_A = T_B \circ f,
\qquad
|\Phi_B(f(x)) - \Phi_A(x)| \le C_\Phi,
\qquad
V_B(f(x)) \le V_A(x) + C_V.
\]


Identities and composition are defined and verified in Lean v1.5.

\section{Collatz as a Canonical Object}

The normalized Collatz macro-step is


\[
T(n) =
\begin{cases}
n/2, & n \text{ even},\

\[4pt]
(3n+1)/2^{v_2(3n+1)}, & n \text{ odd}.
\end{cases}
\]



It defines the concrete discrete dm$^3$-system \texttt{CollatzDm3Candidate} with:


\[
X = \mathbb{N},\quad
T = \texttt{collatz},\quad
\Phi(n) = v_2(n),\quad
V(n) = \log_2(n+1),\quad
\Gamma = \{1,2,4\}.
\]



\subsection{Proved Axioms}

\begin{itemize}
  \item \textbf{operatorDecomposition} (TOGT grammar) --- proved.
  \item \textbf{contactForm} (2-adic dissipation on odd branch) --- proved.
\end{itemize}

\subsection{Analytic Targets (Open)}

\begin{itemize}
  \item \textbf{meanContraction} (3/4 factor above $g_6 = 33$),
  \item \textbf{lyapunovDescent} (average $\Delta V < 0$),
  \item \textbf{hasStructuredCycle} (every orbit reaches $\Gamma$).
\end{itemize}

\section{AXLE Target 5 = Membership Problem}

The Collatz conjecture is equivalent to:



\[
\text{The Collatz map } T \text{ is an object of dm}^3_{\mathrm{disc}}.
\]



In Lean, this is literally the statement:

\begin{verbatim}
theorem collatz_converges (n : ℕ) :
  ∃ k, (Function.iterate collatz k n) ∈ collatzCycle
\end{verbatim}

which reduces to filling the three open fields of
\texttt{CollatzDm3Candidate}.

\section{References}

\begin{itemize}
  \item GCM Manifesto \S2.2 (continuous dm$^3$).
  \item TOGT Report, pp.~108--109 (Collatz as corollary of crystal geometry).
  \item \texttt{DiscreteDm3.lean} v1.5 (formal companion, AXLE repository).
\end{itemize}

\end{document}
